% -----------------------------------------------------------------------------
% December 2006 production study
% -----------------------------------------------------------------------------
% This file defines the observation-facing geospace research experiment.  It is
% deliberately separate from validation.tex: C9 and C10 establish whether the
% model reproduces published/archived cutoff products, whereas this section
% pre-registers the additional morphology, timing, and hysteresis analyses that
% will be populated after the frozen production run.

\section{December 2006 Research Experiment: Storm-Time Cutoff Erosion and Memory}
\label{sec:experiment}

\subsection{Scientific Objective and Falsifiable Hypotheses}

The December 2006 solar energetic particle (SEP) event and geomagnetic storm
provide a particularly useful natural experiment because energetic protons
were present while the magnetosphere underwent rapid compression, sustained
southward interplanetary magnetic field (IMF), main-phase erosion, and
recovery. PAMELA measured a rigidity-dependent equatorward displacement of the
SEP cutoff, reaching approximately $7^{\circ}$ for the lowest reported
rigidities, on its approximately 94-min orbit cadence
\cite{Adriani-2016-SW}. Contemporaneous NOAA POES/MetOp measurements provide a
complementary, multi-spacecraft view of lower-rigidity proton boundaries as a
function of magnetic local time (MLT) \cite{Dmitriev-2010-JGR}. The event
therefore permits a test of absolute cutoff latitude, spatial morphology, and
time dependence using one internally consistent IGRF+TS05 trajectory model.

The primary scientific objective is to determine how the rigidity-dependent
geomagnetic shielding boundary changes during storm development and whether
its recovery is described by the instantaneous geomagnetic disturbance alone
or retains measurable dependence on prior solar-wind driving. Four hypotheses
are fixed before inspecting the production results:

\begin{enumerate}
\item The storm-time equatorward displacement is larger and persists longer at
lower rigidity than at higher rigidity.
\item Boundary motion is not longitudinally uniform: its mean latitude,
first-harmonic MLT amplitude, and phase evolve differently during compression,
main phase, and recovery.
\item Dynamic pressure primarily controls the prompt compression response,
whereas southward IMF, SYM-H, and the TS05 history variables $W_1$--$W_6$
better describe the delayed erosion and recovery.
\item Main- and recovery-phase boundaries differ at matched SYM-H, constituting
hysteresis. If the TS05 history variables contain relevant storm memory, they
should reduce this phase-dependent residual relative to an instantaneous-index
description.
\end{enumerate}

These are model--data and numerical sensitivity hypotheses, not claims that a
single empirical-field parameter is a unique physical cause. Solar-wind and
geomagnetic drivers covary strongly, and TS05 is an empirical storm-time field
model rather than a self-consistent plasma simulation \cite{Tsyganenko-2005-JGR}.
The analysis therefore emphasizes prediction, timing, and controlled
counterfactual calculations; causal language is reserved for mechanisms that
cannot be distinguished by the present design.

\subsection{Event Windows and Predefined Landmarks}

The context interval is 12--18 December 2006. The production interval is
14 December 00:00 UT through 17 December 00:00 UT, matching the complete
five-minute TS05 driver used by the C9 validation. A 15-min model cadence is
used throughout the production interval, with five-minute calculations added
around rapid changes in dynamic pressure, IMF $B_z$, and SYM-H. The cadence is
chosen before examining modeled cutoff motion.

Event landmarks are selected algorithmically from the quality-controlled
driver rather than moved to coincide with a modeled extremum. The quiet
reference $t_q$ is the last complete six-hour interval before the storm for
which median $|\mathrm{SYM\mbox{-}H}|\leq20$ nT, the absolute Theil--Sen SYM-H
trend is no greater than 1 nT h$^{-1}$, and no required TS05 input is missing.
The compression time $t_c$ is the epoch of the largest positive five-minute
change in $\ln P_{\mathrm{dyn}}$ during the storm arrival window. The
main-phase time $t_m$ is the
minimum SYM-H. The maximum-erosion time $t_e(R)$ is determined separately from
the modeled mean cutoff at each rigidity and is reported relative to $t_c$ and
$t_m$. Recovery samples are retained only after $t_m$. Matched main/recovery
pairs are then formed using the driver tolerances specified below. This
separation prevents the observational landmarks and the modeled response times
from being defined circularly.

\begin{table*}
\caption{Predefined Time Sampling and Its Scientific Role}
\label{tab:exp-time-sampling}
\centering
\small
\begin{tabular}{p{0.20\textwidth}p{0.21\textwidth}p{0.49\textwidth}}
\toprule
Product & Sampling & Purpose \\
\midrule
Event context & 12--18 December; 5-min drivers & Establish upstream conditions, disturbance indices, data gaps, and event phases. \\
Production survey & 14 December 00:00--17 December 00:00 UT; 15 min & Resolve rigidity-dependent erosion, recovery, MLT morphology, and accessible area. \\
Rapid-transition subset & 5 min around $t_c$, strong southward-$B_z$ changes, $t_m$, and fastest modeled motion & Bound cadence error and quantify prompt versus delayed response. \\
PAMELA comparison & Published orbit midpoints; optionally a five-minute average over each approximately 94-min accumulation interval & Reproduce the C9 observation operator without treating an orbit-average as an instantaneous measurement. \\
POES/MetOp comparison & Native 16-s samples grouped in two-hour windows and three-hour MLT sectors & Reproduce the C10 constellation operator and estimate MLT-resolved boundaries. \\
\bottomrule
\end{tabular}
\end{table*}

\subsection{Data Sources, Provenance, and Processing}

\subsubsection{Solar-Wind and Geomagnetic Drivers}

Near-Earth solar-wind, IMF, and geomagnetic indices are taken from the NASA
Space Physics Data Facility (SPDF) OMNI high-resolution data service
\cite{King-2005-JGR,NASA-SPDF-OMNI-2026}. The production TS05 input is the
official five-minute 2006 OMNI file distributed with the TS05 history
variables, \texttt{2006\_OMNI\_5m\_with\_TS05\_variables}, from the model
author's archive \cite{Tsyganenko-TS05-Drivers-2006}. Each record supplies the
epoch, GSM IMF $B_y$ and $B_z$, solar-wind speed and density, dynamic pressure,
DST or SYM-H-compatible storm index, dipole tilt, and $W_1$--$W_6$. TS05 uses
dynamic pressure, disturbance level, IMF, dipole tilt, and the six history
variables to represent storm-time current systems \cite{Tsyganenko-2005-JGR}.

The preprocessing archive will retain the unmodified downloaded file, source
URL, retrieval date, checksum, parser version, and a normalized table used by
AMPS. Automated checks require monotonic five-minute epochs, documented units,
finite required values, no duplicate epochs, physically plausible ranges, and
an explicit gap mask. No interpolation is allowed across a gap longer than
15 min. Shorter gaps are linearly interpolated for instantaneous quantities;
the interpolation flag is retained. Dynamic pressure is recomputed from the
available density and speed as a consistency check but the distributed TS05
value remains authoritative unless a documented file defect is found. The
exact row or interpolating pair selected for every AMPS snapshot is written to
the run manifest. These checks extend the C15 driver-interpolation contract.

\subsubsection{PAMELA Rigidity-Dependent Cutoffs}

The PAMELA comparison uses the numerical supporting information to
\citet{Adriani-2016-SW}, specifically Table S1 rather than values digitized
from a figure. The table contains 37 orbit midpoints and seven rigidity bins,
yielding 259 time--rigidity cells, of which 258 contain reported cutoffs. The
geometric-mean bin rigidities are 0.424, 0.498, 0.588, 0.693, 0.817, 0.962,
and 1.131 GV. PAMELA defined the observed cutoff latitude as the latitude at
which the measured proton flux reached one half of its mean high-latitude
value above $65^{\circ}$ AACGM latitude. Published asymmetric uncertainties
are preserved. The primary analysis uses the published orbit midpoint and
cutoff directly, because the public table is an orbit-level product rather
than event data. A temporal-sampling sensitivity calculation averages model
transmission over five-minute snapshots spanning the nominal 94-min orbit
window; it is reported separately and is not represented as reconstruction of
the unpublished event selection.

The validation-equivalent model samples a 475-km WGS-84 geodetic shell, traces
vertical protons, longitude-averages the resolved access state separately in
each hemisphere, and extracts a weighted-isotonic 50\% transmission crossing.
The north and south absolute cutoff latitudes are combined using their median,
matching the C9 contract. This shell approximation is an explicit assumption:
it removes orbit-location details not present in Table S1 and tests the global
cutoff response rather than a spacecraft-by-spacecraft ephemeris reconstruction.

\subsubsection{NOAA POES/MetOp Proton Boundaries}

The local-time analysis uses NOAA/NCEI Level-2 SEM-2 16-s data from NOAA-15,
NOAA-16, NOAA-17, NOAA-18, and MetOp-A (MetOp-02), obtained from the public
POES/MetOp Space Environment Monitor archive \cite{NOAA-NCEI-POES-SEM2}.
Required fields are UTC, geodetic latitude, longitude, altitude, $L$, MLT, and
the MEPED omnidirectional proton variables \texttt{mepomp6} through
\texttt{mepomp9}. The archived daily files, retrieval URLs, checksums, and the
NOAA processing-version identifier are recorded. Samples outside 14 December
00:00 UT--16 December 12:00 UT, samples with nonzero fatal quality flags, and
nonfinite or nonpositive fluxes are excluded with reason codes. Positions are
converted to AACGM latitude and MLT using one frozen coordinate-library
version.

The four integral thresholds are represented by the lower-threshold proton
rigidities 0.174013525, 0.262395866, 0.369131538, and 0.531334344 GV for P6,
P7, P8, and P9, respectively. P6 and P7 are primary. P8 and P9 are diagnostics
because their subcommutated sampling, broad response, and spectral dependence
make a sharp lower-threshold interpretation less secure
\cite{Evans-2000-NOAA,Machol-2012-NOAA,Green-2013-NOAA,Yando-2011-JGR}.
Additional limitations include proton/electron contamination, radiation
damage and threshold drift, broad angular response, and cross-spacecraft
calibration differences \cite{Asikainen-2011-JASTP,Sandanger-2015-JGR,Peck-2015-JGR}.
Consequently, the model is not empirically adjusted to make these channels
agree, and the primary claims do not rely on P8/P9.

For each poleward pass, the observed normalized transmission is
\begin{linenomath*}
\begin{equation}
 T_{\mathrm{obs}}(\lambda)=
 \frac{F(\lambda)-F_{\mathrm{bg}}}
      {F_{\mathrm{pol}}-F_{\mathrm{bg}}},
 \label{eq:exp-poes-transmission}
\end{equation}
\end{linenomath*}
where $F_{\mathrm{pol}}$ is the median flux poleward of $75^{\circ}$ AACGM and
$F_{\mathrm{bg}}$ is the median within $5^{\circ}$ of the low-latitude edge of
the retained leg. A weighted isotonic fit gives $T_{25}$, $T_{50}$, and
$T_{75}$. Eligible crossings require the sample counts, polar-to-background
contrast, bracket, width, and isotonic-residual criteria specified by C10.
Eligible legs are grouped into two-hour windows centered hourly and into eight
three-hour MLT sectors. The corresponding model calculation uses an 850-km
shell and the identical AACGM/MLT binning and crossing estimator.

\begin{table*}
\caption{Observation and Driver Data Used in the December 2006 Study}
\label{tab:exp-data}
\centering
\small
\begin{tabular}{p{0.18\textwidth}p{0.21\textwidth}p{0.24\textwidth}p{0.27\textwidth}}
\toprule
Data set & Source & Variables retained & Role and principal limitation \\
\midrule
Five-minute OMNI/TS05 driver & NASA SPDF OMNI and official TS05 driver archive & $P_{\mathrm{dyn}}$, GSM $B_y,B_z$, speed, density, DST/SYM-H, tilt, $W_1$--$W_6$ & Drives IGRF+TS05; empirical inputs are correlated and carry upstream timing and interpolation uncertainty. \\
PAMELA Table S1 & Supporting information to \citet{Adriani-2016-SW} & Orbit midpoint, rigidity bin, absolute cutoff latitude, asymmetric uncertainty & Primary rigidity-dependent absolute validation; orbit-averaged product lacks event-level selection and full ephemeris. \\
NOAA POES/MetOp SEM-2 Level 2 & NOAA/NCEI archive & 16-s position, MLT, quality, P6--P9 omnidirectional proton channels & Primary P6/P7 MLT-resolved validation; broad response, aging, contamination, and threshold-proxy assumptions. \\
AMPS run manifest & Generated with every production case & Input checksums, exact driver rows, command, revision, numerical controls, state counts & Makes every modeled point reproducible and prevents silent mixing of configurations. \\
\bottomrule
\end{tabular}
\end{table*}

\subsection{Production Model and Numerical Design}

The production magnetic field is IGRF plus TS05. Trajectory tracing in
realistic, time-indexed geomagnetic fields is the established basis for
storm-time cutoff calculations \cite{Cooke-1991-INC,Smart-2005-ASR,
Kress-2010-SW,Kress-2015-JGR}. No alternative external-field
model enters the primary science result; alternative fields, if calculated,
are labeled structural-sensitivity experiments. The field is frozen during
each backtrace, and the electric field is set to zero, so the result is a
quasi-static rigidity access boundary rather than a time-dependent
electromagnetic transport solution. Relativistic full-orbit protons are
backtraced with the RK4 mover. Boris is used at selected quiet, compression,
main-phase, and recovery epochs to test mover sensitivity. Mode3D mesh field
evaluation supplies the production time series, while GRIDLESS field
evaluation and mesh refinement provide parity and interpolation checks.

The trajectory classifier returns three states,
\begin{linenomath*}
\begin{equation}
 A(R,\Omega,\mathbf{x},t) \in \{0,1,2\},
 \label{eq:exp-access-state}
\end{equation}
\end{linenomath*}
for forbidden, allowed, and unresolved trajectories. Only states 0 and 1 enter
the resolved transmission,
\begin{linenomath*}
\begin{equation}
 T_{\mathrm{mod}}(R)=
 \frac{\sum_{j:A_j\ne2} w_j A_j}
      {\sum_{j:A_j\ne2} w_j}, \qquad
 f_{\mathrm{unres}}=
 \frac{\sum_{j:A_j=2}w_j}{\sum_j w_j},
 \label{eq:exp-model-transmission}
\end{equation}
\end{linenomath*}
where $j$ indexes longitude or look direction and $w_j$ is its quadrature
weight. Unresolved weight is never silently converted to allowed or forbidden.
A reported boundary must satisfy the C9 or C10 resolved-coverage and bracketing
requirements; sensitivity runs lengthen the trace limit to distinguish a
physical penumbra from insufficient integration.

One common spatial product supplies two observation-equivalent reductions and
the research analysis. It contains the C9 and C10 shell altitudes, vertical
incidence, exact rigidity samples, exact observation midpoints, and all inputs
needed by their separate time-aggregation and $T_{50}$ operators. The research
product uses both 475- and 850-km shells, a $15^{\circ}$ longitude by
$2^{\circ}$ latitude grid from $35^{\circ}$ to $85^{\circ}$ in both
hemispheres, and a common rigidity grid from 0.15 to 1.25 GV. Exact PAMELA and
POES/MetOp rigidities are inserted into that grid. Direct access is evaluated
at each rigidity; cutoff boundaries are not inferred by interpolating between
widely separated energy channels. The primary maps use vertical incidence so
that they remain commensurate with the cutoff-latitude validation. Directional
and aperture-folded access are reserved for a labeled secondary analysis and
are not mixed with vertical cutoffs.

The two shells are evaluated in one Mode3D process. Restartable groups of up to
16 irregular epochs use \texttt{TEMPORAL\_MODE SNAPSHOT\_LIST}: AMR topology,
distributed block allocation, and static data layout are retained, but the
five-minute TS05 driver is interpolated and the complete IGRF+TS05 field is
refilled at every epoch. The optimized result is required to reproduce a
standalone one-epoch/one-shell calculation with complete access-grid closure,
identical resolved access states and unresolved fractions, and derived
$T_{50}$ boundaries equal to numerical reporting precision. Thus mesh reuse
changes process lifecycle but not the physical snapshot definition.

\subsection{Observation Operators and Model--Data Comparisons}

For an observation $i$ and its strictly paired model value, the latitude
residual is
\begin{linenomath*}
\begin{equation}
 \epsilon_i=\Lambda_{c,i}^{\mathrm{mod}}-
             \Lambda_{c,i}^{\mathrm{obs}}.
 \label{eq:exp-residual}
\end{equation}
\end{linenomath*}
Positive $\epsilon_i$ denotes excessive modeled shielding (a boundary that is
too poleward), and negative $\epsilon_i$ denotes excessive modeled access. The
reported aggregate measures are bias, median absolute error, root-mean-square
error, Pearson and rank correlation, and the timing and magnitude errors of
the maximum equatorward displacement. PAMELA metrics are reported overall and
by rigidity. POES/MetOp metrics are reported for primary P6/P7 cells overall
and by MLT, hemisphere, spacecraft, and storm phase; P8/P9 results remain
diagnostic. All denominators use paired eligible cells, and tables show the
eligible, rejected, missing, and unresolved counts.

PAMELA and POES/MetOp are not treated as interchangeable measurements. Their
altitudes, rigidities, directions, sampling, and cutoff definitions differ.
Cross-instrument closure is tested only after imposing PAMELA's temporal and
polar-pass sampling on the POES/MetOp constellation and applying the
corresponding operator to AMPS. This nested comparison asks whether the model
explains the expected altitude, rigidity, MLT, and sampling differences; it
does not force either instrument onto the other's native measurement.

\begin{table*}
\caption{Primary Quantities Compared With Observations}
\label{tab:exp-comparisons}
\centering
\small
\begin{tabular}{p{0.20\textwidth}p{0.25\textwidth}p{0.25\textwidth}p{0.20\textwidth}}
\toprule
Quantity & Observation & AMPS analogue & Scientific use \\
\midrule
PAMELA cutoff latitude & Published 50\% polar-flux latitude in seven rigidity bins & Longitude-averaged resolved vertical-access $T_{50}$ at 475 km & Absolute rigidity-dependent validation and storm displacement. \\
POES/MetOp P6/P7 boundary & Isotonic $T_{50}$ from normalized leg profiles & Identically binned vertical-access $T_{50}$ at 850 km & Low-rigidity MLT, hemisphere, and time validation. \\
Transition width & $T_{75}-T_{25}$ from eligible NOAA legs & Modeled $T_{75}-T_{25}$ & Penumbra sharpness; secondary because detector response broadens the observed transition. \\
Boundary timing & Epoch of largest observed equatorward displacement & Epoch of modeled minimum cutoff latitude & Prompt-response and recovery timing. \\
Cross-instrument closure & Sampling-matched PAMELA and POES/MetOp boundaries & Same two observation operators applied to one field solution & Consistency across altitude, rigidity, and sampling. \\
\bottomrule
\end{tabular}
\end{table*}

\subsection{Cutoff Morphology and Accessible-Area Diagnostics}

At each epoch, rigidity, altitude, and hemisphere, the MLT-resolved boundary is
fit first with a mean and first harmonic,
\begin{linenomath*}
\begin{equation}
 \Lambda_c(\mathrm{MLT},t,R)=\Lambda_0(t,R)+A_1(t,R)
 \cos\left[\frac{2\pi}{24}\left(\mathrm{MLT}-\phi_1(t,R)\right)\right].
 \label{eq:exp-harmonic}
\end{equation}
\end{linenomath*}
The implemented reduction uses ordinary least squares on finite resolved MLT
cells and archives the valid/requested cell counts and RMS residual. A second
fit adds $C_2\cos(2\theta)+S_2\sin(2\theta)$, where
$\theta=2\pi\,\mathrm{MLT}/24$, and reports
$A_2=(C_2^2+S_2^2)^{1/2}$ and its 12-h phase. This semidiurnal term is a
descriptive day--night versus dawn--dusk deformation metric; the original cell
table and both fit residuals remain authoritative. The storm perturbation is
referenced to the
quiet median,
\begin{linenomath*}
\begin{equation}
 \Delta\Lambda_c(t,R,\mathrm{MLT})=
 \Lambda_c(t,R,\mathrm{MLT})-
 \widetilde{\Lambda}_{c,q}(R,\mathrm{MLT}),
 \label{eq:exp-cutoff-depression}
\end{equation}
\end{linenomath*}
so negative values are equatorward erosion. Boundary speed is evaluated with
centered finite differences except at gaps and interval endpoints.

The implemented boundary-derived accessible fraction within the analyzed
absolute-latitude band $[\Lambda_{\min},\Lambda_{\max}]$ is
\begin{linenomath*}
\begin{equation}
 f_{\mathrm{acc}}(R,t)=\left\langle
 \frac{\sin\Lambda_{\max}-\sin\Lambda_c(\mathrm{MLT},t,R)}
 {\sin\Lambda_{\max}-\sin\Lambda_{\min}}
 \right\rangle_{\mathrm{MLT}},\qquad
 S_{\mathrm{eq}}=2\pi(R_E+h)^2
 (\sin\Lambda_{\max}-\sin\Lambda_{\min})f_{\mathrm{acc}},
 \label{eq:exp-accessible-area}
\end{equation}
\end{linenomath*}
where boundaries are clamped to the analyzed band before integration and the
average is over valid MLT sectors. $S_{\mathrm{eq}}$ is explicitly a
spherical-shell-equivalent comparison metric; AACGM is not assumed to be an
equal-area coordinate system. Reporting both latitude and accessible fraction
distinguishes a nearly uniform equatorward shift from a strongly asymmetric
opening confined to a limited range of MLT.

The same exact-rigidity classifications are inverted independently at every
geographic shell cell to form spatial cutoff-rigidity maps. Unresolved states
are omitted, and the remaining binary access sequence $A(R_i)\in\{0,1\}$ is
projected onto a nondecreasing sequence $\widehat{A}(R_i)$ using the
equal-weight pool-adjacent-violators algorithm. A numerical map value is
reported only when
\begin{linenomath*}
\begin{equation}
 \widehat{A}(R_{\min})<0.5<\widehat{A}(R_{\max}), \qquad
 R_{50}:\ \widehat{A}(R_{50})=0.5,
 \label{eq:exp-spatial-r50}
\end{equation}
\end{linenomath*}
with linear interpolation between adjacent fitted samples; an exact 0.5
plateau is represented by the midpoint of its rigidity interval. All-allowed
and all-forbidden profiles are retained as $R_{50}<R_{\min}$ and
$R_{50}>R_{\max}$ censoring states rather than assigned endpoint values.
Mixed profiles without an endpoint bracket and incomplete profiles remain
non-numerical. Bracket width, resolved fraction, and raw nonmonotonic transition
counts accompany every map cell.

For a cell $j$ with a bracketed precompression quiet median $R_{50,q,j}$, the
storm-time decrease is
\begin{linenomath*}
\begin{equation}
 D_j(t)=\max[0,R_{50,q,j}-R_{50,j}(t)].
 \label{eq:exp-spatial-cutoff-decrease}
\end{equation}
\end{linenomath*}
The maximum of $D_j(t)$ over the event identifies where and when shielding loss
is greatest. If the event profile becomes all-allowed, the recorded value
$R_{50,q,j}-R_{\min}$ is a conservative lower bound and is flagged as censored;
it is excluded from exact mean-change statistics. Shell-level evolution uses
$\cos(\lambda_{\mathrm{GEO}})$ area weights and reports exact mean change,
conservative median/90th-percentile/maximum decrease, and the weighted area
fraction for which $D_j\geq0.05$~GV.

\subsection{Lag, Driver Attribution, and Hysteresis}

The candidate drivers are $P_{\mathrm{dyn}}$, IMF $B_z$, the time integral of
southward IMF, SYM-H, and $W_1$--$W_6$. For a cutoff diagnostic $y(t)$ and a
standardized driver $x(t)$, lagged association is summarized by
\begin{linenomath*}
\begin{equation}
 C_{xy}(\tau)=
 \frac{\sum_{t\in\mathcal{V}_{\tau}}[x(t-\tau)-\bar{x}_{\tau}]
 [y(t)-\bar{y}_{\tau}]}
 {\left\{\sum[x(t-\tau)-\bar{x}_{\tau}]^2
 \sum[y(t)-\bar{y}_{\tau}]^2\right\}^{1/2}},
 \label{eq:exp-lag-correlation}
\end{equation}
\end{linenomath*}
where $\mathcal{V}_{\tau}$ contains only simultaneous valid samples. Lags are
searched over the preregistered interval $-3$ to $+6$ h, with positive $\tau$
meaning that cutoff response follows the driver. Confidence intervals use
phase-preserving block bootstrap resamples; the effective number of
independent samples, not the raw five-minute count, controls significance.
Runs with fewer than 24 paired model epochs retain lag curves for software and
exploratory diagnostics but omit bootstrap intervals and are labeled
\texttt{DIAGNOSTIC\_ONLY}. The same minimum is applied to matched-driver
hysteresis inference. Recovery half-times and e-folding times additionally
require at least six post-peak samples and observed crossings of both levels.

Predictive attribution uses a regularized distributed-lag model,
\begin{linenomath*}
\begin{equation}
 y(t,R)=\beta_0(R)+\sum_m\sum_{k\in\mathcal{K}}
 \beta_{mk}(R)x_m(t-k)+\gamma(R)\,\mathrm{MLT}+\varepsilon(t,R),
 \label{eq:exp-distributed-lag}
\end{equation}
\end{linenomath*}
with contiguous storm-phase blocks for cross-validation. Ridge or elastic-net
regularization is selected within the training blocks and performance is
reported on held-out blocks. The baseline contains SYM-H and dynamic pressure;
the history model adds $B_z$ and $W_1$--$W_6$. Improvement is measured by
held-out RMSE and correlation. Coefficient magnitude is not interpreted as a
causal partition because the predictors are correlated and the $W_i$ are
themselves derived driver histories.

Hysteresis is tested without fitting a visual loop. Each main-phase sample is
matched to a recovery-phase sample at the same rigidity, hemisphere, and MLT,
with $|\Delta\mathrm{SYM\mbox{-}H}|\leq10$ nT. A stricter subset additionally
requires dynamic pressure within 20\% and $|\Delta B_z|\leq2$ nT. The matched
phase contrast is
\begin{linenomath*}
\begin{equation}
 \Delta\Lambda_{\mathrm{hys}}(R,\mathrm{MLT})=
 \operatorname{median}_{p}
 \left[\Lambda_{c,p}^{\mathrm{rec}}-
       \Lambda_{c,p}^{\mathrm{main}}\right],
 \label{eq:exp-hysteresis}
\end{equation}
\end{linenomath*}
and the signed loop area $\oint\Lambda_0\,\mathrm{d}(\mathrm{SYM\mbox{-}H})$
is a secondary descriptor. Bootstrap confidence intervals resample matched
pairs in contiguous time blocks. The history hypothesis is supported only if
the contrast is resolved numerically, has a confidence interval excluding
zero in the preregistered primary rigidity set, and is reduced in held-out
prediction after $W_1$--$W_6$ are included.

Three TS05 sensitivity experiments help interpret this result: the full
time-varying driver; a calculation in which the $W_i$ are held at quiet values
while instantaneous pressure, IMF, and disturbance inputs vary; and a
calculation in which instantaneous inputs are held while the $W_i$ evolve.
These are controlled numerical perturbations of an empirical model, not
self-consistent alternative magnetospheres. Their purpose is to identify which
parts of the TS05 parameterization generate modeled memory, not to prove that a
particular magnetospheric current system caused the observations.

\subsection{Numerical Convergence, Uncertainty, and Decision Rules}

The numerical ensemble varies spatial mesh resolution, trajectory time step,
trace duration, longitude/MLT resolution, rigidity spacing, mover, and
Mode3D versus GRIDLESS field evaluation. The production result is accepted
only if its key diagnostics are stable relative to the observational
uncertainty and the unresolved fraction remains below the C9/C10 limits.
Scheduler, MPI-rank, and thread-count invariance are checked by C13 and C18 and
are not counted as independent physical ensembles. C15 verifies that the
selected TS05 epoch and driver interpolation are reproducible.

Observational uncertainty is kept separate from numerical spread. PAMELA uses
the published asymmetric intervals and orbit-level sampling uncertainty.
POES/MetOp uses block variation among legs and spacecraft, together with
sensitivity to plateau/background windows, quality screening, and channel
interpretation. AACGM transformation version, altitude-shell approximation,
and TS05 structural limitations are reported as model-form assumptions rather
than folded into a single Gaussian error bar. Confidence intervals are based
on contiguous block bootstrap resampling. Rigidity- or MLT-resolved multiple
comparisons use a false-discovery-rate correction, while the small set of
primary global metrics is reported with unadjusted confidence intervals.

The established C9 acceptance targets (valid fraction at least 0.85, RMSE no
greater than $3^{\circ}$, absolute bias no greater than $2^{\circ}$, and
correlation at least 0.80) remain the formal PAMELA validation gates. C10's
cell-quality and coverage rules remain the POES/MetOp gates. The new science
conclusions are not declared successful merely because those validation gates
pass: the effect size, confidence interval, numerical stability, and dependence
on sampling and model assumptions are reported independently.

\begin{table*}
\caption{Planned Analysis Products and Predefined Interpretation}
\label{tab:exp-products}
\centering
\small
\begin{tabular}{p{0.21\textwidth}p{0.34\textwidth}p{0.35\textwidth}}
\toprule
Product & Reported quantities & Decision or interpretation \\
\midrule
Absolute validation & Bias, MAE, RMSE, correlation, valid/rejected counts, maximum-displacement timing & Apply unchanged C9/C10 gates; report all paired residuals. \\
Rigidity dependence & $\Delta\Lambda_c$, $S_{\mathrm{acc}}$, time of maximum erosion versus $R$ & Test ordered erosion and recovery with uncertainty, not only a selected low/high pair. \\
MLT morphology & $\Lambda_0$, $A_1$, $\phi_1$, optional cross-validated second harmonic & Identify global shift versus localized access and compare with POES/MetOp sectors. \\
Lag response & Peak $C_{xy}(\tau)$, lag interval, held-out prediction skill & Describe predictive timing; do not infer unique causation from correlated drivers. \\
Hysteresis & Matched-phase contrast, loop area, rigidity/MLT dependence & Require resolved, numerically stable contrast and report unmatched/rejected samples. \\
TS05 memory sensitivity & Difference between full, history-frozen, and instantaneous-frozen experiments & Attribute behavior to the empirical parameterization only; compare with observed phase residuals. \\
Robustness & Mesh, mover, timestep, trace-limit, cadence, shell, and observation-operator spread & A claimed feature must exceed numerical spread and persist under the preregistered sensitivities. \\
\bottomrule
\end{tabular}
\end{table*}

\subsection{Results-Section Contract}

The production archive will populate the following result sequence without
changing definitions after the values are known. First, a driver panel will
show $P_{\mathrm{dyn}}$, IMF $B_z$, SYM-H, and $W_1$--$W_6$, with the objective
event landmarks and all data gaps. Second, paired PAMELA and POES/MetOp panels
will establish absolute model skill and residual structure. Third, maps at
$t_q$, $t_c$, $t_m$ or $t_e$, and a matched recovery time will show cutoff
erosion and accessible area at selected rigidities. Fourth, time series of
$\Lambda_0$, $A_1$, $\phi_1$, $S_{\mathrm{acc}}$, and unresolved fraction will
quantify morphology and numerical validity. Finally, lag curves, matched-phase
hysteresis, and the three TS05 sensitivity experiments will test the memory
hypothesis. Machine-readable tables will accompany every plotted summary.

This ordering makes negative results interpretable. A good absolute cutoff
comparison with no resolved hysteresis would support AMPS as a storm-time
shielding model while placing an upper bound on memory detectable in this
event. Conversely, a modeled hysteresis signal without corresponding
observation support, or one that disappears under numerical refinement, would
identify a limitation of the TS05-based configuration rather than a discovery.
The placeholders in Section~\ref{sec:results} are therefore a reporting
contract: they will be replaced only by values from the frozen archive,
including uncertainty, sample count, and sensitivity range.
